\documentclass{article}
\usepackage{iclr2026_conference,times}
\input{math_commands.tex}
\usepackage{booktabs}
\usepackage{hyperref}
\usepackage{url}

\title{Interim Report: Module-Aware Low-Bit Quantization for DINO-WM Planning}
\author{}

\begin{document}
\maketitle
\lhead{Under review as a workshop paper at ICLR 2026 (ES-Reasoning)}

\begin{abstract}
This short report summarizes in-progress results from an active quantization study on DINO-WM planning for Wall.
The main research question is whether \emph{where} precision is allocated matters as much as \emph{how many} bits are used.
At the time of writing, paired-target generation completed and the strict 4-hour core experiment block is still running.
One finalized baseline metric is available (FP16, budget bA, seed 1): success rate 0.50 over 10 evaluations.
We document current artifacts, sanity checks, and the exact experimental schedule that remains to finish the main claim comparisons (mixed INT4 vs uniform INT4, with FP16/INT8 anchors).
\end{abstract}

\section{Goal and Scope}
Our goal is to validate two claims for world-model planning:
\begin{enumerate}
  \item Near 4-bit precision, allocation structure matters: mixed INT4 (encoder preserved, predictor quantized) should outperform uniform INT4.
  \item Mixed INT4 should remain reasonably close to stronger anchors (INT8 and FP16) under matched planning budgets.
\end{enumerate}

This report is intentionally \emph{interim}: we only summarize completed artifacts and avoid overclaiming unfinished comparisons.

\section{Current Experimental State}
\paragraph{Environment and model.}
All runs use DINO-WM Wall with checkpoint \texttt{wall\_single}, running on Apple Silicon (MPS backend).

\paragraph{Artifacts completed so far.}
\begin{itemize}
  \item Paired-target files generated: 8 files across budgets \texttt{bA}, \texttt{bB}, \texttt{bC}.
  \item Paired-target manifest written:
  \texttt{results/paired\_targets/m3\_main2exp\_v2/paired\_targets\_manifest.json}.
  \item Finalized planning metrics available: 1 file so far
  (\texttt{fp16/budget\_bA/seed\_1/metrics.json}).
\end{itemize}

\paragraph{Strict 4-hour run plan (active).}
The active schedule keeps only the core claim blocks:
\begin{itemize}
  \item Core frontier on \texttt{bA,bB}, seeds \{1,2\}: \texttt{fp16}, \texttt{uniform\_int8}, \texttt{mixed\_int8}, \texttt{uniform\_int4}, \texttt{mixed\_int4}.
  \item Module-location check on \texttt{bA, seed1}: \texttt{encfp16\_50} vs \texttt{encheadfp16\_50}.
\end{itemize}
This produces 22 required run cells in the strict plan.

\section{Interim Results}
\begin{table}[t]
\centering
\small
\begin{tabular}{lcc}
\toprule
Metric (available now) & Value & Source \\
\midrule
Paired target files & 8 & paired\_targets tree \\
Finalized metrics files & 1 & wall\_m3\_main2\_expansion tree \\
FP16 success (bA, seed1) & 0.50 (5/10) & metrics.json \\
FP16 avg plan time (bA, seed1) & 252.7 s & metrics.json \\
FP16 model size & 205.0 MB & metrics.json \\
\bottomrule
\end{tabular}
\caption{Interim results available at manuscript time.}
\label{tab:interim}
\end{table}

The single finalized anchor point (FP16, bA, seed1) is non-trivial (0.50 success), which is a useful sanity baseline for the pending low-bit comparisons.
However, no mixed-vs-uniform INT4 paired deltas are finalized yet; those are still in progress.

\section{Sanity Checks}
We ran three sanity checks on current artifacts:
\begin{itemize}
  \item \textbf{Write integrity}: output directories, manifests, and per-run logs are present and internally consistent.
  \item \textbf{Pipeline progression}: logs show transition from paired-target generation to \texttt{run\_wall\_grid}.
  \item \textbf{Resource headroom}: disk free space remains above critical threshold during current strict run.
\end{itemize}

One cautionary observation is that paired-target manifest entries report generator success rates of 0.0.
Because target files are still produced and downstream runs proceed, this may reflect target-generation dynamics rather than a fatal pipeline issue, but it should be audited before final claims.

\section{What Remains for Final Claims}
The following comparisons are still required to support the main paper claim set:
\begin{itemize}
  \item \textbf{Primary}: \texttt{mixed\_int4} vs \texttt{uniform\_int4} on \texttt{bA,bB}, seeds \{1,2\}.
  \item \textbf{Anchor distance}: compare mixed INT4 to FP16 and INT8 variants under the same budgets.
  \item \textbf{Where-bits test}: \texttt{encfp16\_50} vs \texttt{encheadfp16\_50} on \texttt{bA, seed1}.
\end{itemize}
Once these finish, we will run aggregation and paired-delta scripts and finalize statistical conclusions.

\section{Conclusion}
The experiment infrastructure is functional and currently producing results under a strict runtime budget.
At this interim point, we have one validated FP16 anchor and complete paired-target artifacts, but not enough finalized low-bit metrics to claim mixed INT4 superiority yet.
The remaining strict schedule is targeted specifically at that claim and should produce a decisive mixed-vs-uniform INT4 comparison with FP16/INT8 context.

\bibliography{references}
\bibliographystyle{iclr2026_conference}

\end{document}
